\documentclass[11pt]{article}
\usepackage[margin=2.2cm]{geometry}
\usepackage{amsmath, amssymb}
\usepackage{parskip}
\usepackage{enumitem}
\usepackage{xcolor}
\usepackage{listings}
\usepackage[colorlinks=true, linkcolor=blue!50!black, urlcolor=blue!50!black]{hyperref}

\lstset{
  basicstyle=\ttfamily\small,
  backgroundcolor=\color{black!4},
  frame=single, framerule=0.3pt, rulecolor=\color{black!30},
  columns=fullflexible, keepspaces=true, breaklines=true,
  aboveskip=6pt, belowskip=6pt,
}

\newcommand{\file}[1]{\texttt{#1}}
\newcommand{\code}[1]{\texttt{#1}}

% One header per problem, mimicking the real online sheets.
\newcommand{\problemheader}[3]{%
  \begin{center}
    {\large January 2026 CSE 220}\\
    {\large Practice Online on DFT and FFT}\\[2pt]
    {\Large\bfseries #1}\\[2pt]
    Subsections: #2\\
    Total Marks: 10\\
    Total Time: 40 Minutes\\[2pt]
    {\small Template: \file{#3}}
  \end{center}
  \vspace{4pt}
}

\begin{document}

%=====================================================================
\begin{center}
  {\Large\bfseries Practice Onlines: DFT \& FFT}\\[4pt]
  {\large Ten problems in the format of Online 4 (A1/A2, B1/B2)}\\[2pt]
  five on Task A (big-integer multiplication), five on Task B (image convolution)
\end{center}

\subsection*{How to use this set}

\begin{itemize}[nosep]
  \item Templates are in \file{practice/problems/}, full solutions in
        \file{practice/solutions/}. Each template has exactly three \code{TODO}
        blocks and a provided runner with an independent oracle that prints
        \code{verification: MATCH} or \code{MISMATCH}, exactly like the real sheets.
  \item Run a template \emph{unmodified} first and watch which \code{TODO} it
        stops at. Then fill the three blocks. Then run it, and for Task B run it
        again with \code{--color}.
  \item Every file imports your own \file{transforms.py}, \file{bigmul.py},
        \file{image\_conv.py}. The only difference from a real online is the
        two \code{sys.path} lines at the top, so the files can run from the
        \file{practice/} folder instead of being copied next to the offline
        files. The real sheet's command is \code{python3 ...}; on your machine
        it is \code{python ...}.
  \item Outputs land in \file{practice/outputs/<problem>/}. Every report must
        end with \code{verification: MATCH}; an error above $10^{-9}$ means a bug.
  \item Rules as in the real onlines: NumPy array arithmetic and the functions
        the template imports only; \code{numpy.fft}, SciPy transforms and
        library convolution routines are prohibited; Python big integers only
        in the verification line.
  \item Suggested order: fill \code{TODO 3} (dispatch/glue) first, then
        \code{TODO 1} (helper), then \code{TODO 2} (the core). Time yourself:
        40 minutes each.
\end{itemize}

Solutions in plain language begin on page~12.

%=====================================================================
\newpage
\problemheader{Problem A1: Product of Three Integers}{A1, A2}{a1\_triple\_product.py}

The convolution theorem does not care how many factors there are. The digit
array of $a \cdot b \cdot c$ is the linear convolution $a * b * c$, and in the
frequency domain that is the pointwise product of \emph{three} spectra. Three
forward transforms, one triple product, one inverse transform.

You are given two operands from an offline input file and a third reproducible
operand in \file{a1\_triple\_product.py}. Complete only the three blocks marked
``TODO''. Do not modify \file{transforms.py}, \file{bigmul.py}, or any other
offline file.

\textbf{Mathematical model.} Let $a$, $b$, $c$ be little-endian limb arrays of
lengths $n$, $q$, $r$ in base $B$, with spectra $A$, $B$, $C$ at a common length
$N$. Their product has $n + q + r - 2$ coefficients and
\[
  P[k] = A[k]\,B[k]\,C[k], \qquad p = \mathrm{IDFT}\{P\}, \qquad N \ge n + q + r - 2 .
\]
A coefficient of a triple product is a sum of roughly $n^2/2$ products of three
limbs, so it can reach $(n^2/2)(B-1)^3$. With $B = 10^4$ that exceeds the
$2^{53}$ mantissa for the provided inputs; the template therefore fixes
\code{TRIPLE\_BASE\_DIGITS = 2}, i.e.\ $B = 10^2$. Pass that base to both
\code{to\_limbs} and \code{from\_limbs}. Signs stay outside the transform and
multiply.

\textbf{Tasks}
\begin{enumerate}[nosep]
  \item Complete \code{choose\_transform\_length}. Return the minimum linear
        length for three factors, rounded to what the selected engine supports.
        \hfill [2 Marks]
  \item Complete \code{multiply\_three\_transform}. Zero-pad, transform each
        operand once, multiply the three spectra, perform exactly one inverse
        transform, keep the real part, crop, round to \code{int64}.
        \hfill [5 Marks]
  \item Complete \code{multiply\_three}. Convert with the triple base, combine
        the signs, carry with the same base, return the product string.
        \hfill [3 Marks]
\end{enumerate}

\textbf{Running and verification.}
\begin{lstlisting}
python a1_triple_product.py --input inputs/3.txt --third-digits 2000 --engine fft
\end{lstlisting}
The runner writes \file{product.txt} and \file{report.txt} and compares against
\code{int(a) * int(b) * int(c)} --- the only permitted use of Python's big
integers. The report must show \code{verification: MATCH}. Negative operands
(\file{inputs/1.txt}) must also work.

%=====================================================================
\newpage
\problemheader{Problem A2: $a\,b + c\,d$ With One Inverse Transform}{A1, A2}{a2\_sum\_of\_products.py}

The DFT is linear: the transform of a sum is the sum of the transforms. The sum
of two products therefore needs four forward transforms but only \emph{one}
inverse transform --- add the two product spectra first, invert once.

You are given four reproducible non-negative operands in
\file{a2\_sum\_of\_products.py}. Complete only the three blocks marked ``TODO''.
Do not modify any offline file.

\textbf{Mathematical model.} For limb arrays $a, b, c, d$ with spectra
$A, B, C, D$ at a common length $N$,
\[
  S[k] = A[k]\,B[k] + C[k]\,D[k], \qquad s = \mathrm{IDFT}\{S\},
  \qquad N \ge \max(n_a + n_b - 1,\; n_c + n_d - 1).
\]
The coefficients of $s$ are the un-carried digits of $a b + c d$; one carry
sweep turns them into the decimal answer. All operands are non-negative: a
signed sum can produce negative coefficients, which the offline's
\code{from\_limbs} cannot carry, so a negative operand must raise
\code{ValueError}.

\textbf{Tasks}
\begin{enumerate}[nosep]
  \item Complete \code{choose\_transform\_length}. One length must hold the
        longer of the two products; round to what the engine supports.
        \hfill [2 Marks]
  \item Complete \code{sum\_of\_products\_transform}. Pad and transform all
        four arrays, form $AB + CD$, perform exactly one inverse transform,
        keep the real part, crop to the longer linear length, round.
        \hfill [5 Marks]
  \item Complete \code{sum\_of\_products}. Convert, reject negatives, convolve,
        carry, return the decimal string.
        \hfill [3 Marks]
\end{enumerate}

\textbf{Running and verification.}
\begin{lstlisting}
python a2_sum_of_products.py --digits 1500 --seed 11 --engine fft
\end{lstlisting}
The runner compares against \code{int(a)*int(b) + int(c)*int(d)}. The report
must show \code{verification: MATCH}. Two inverse transforms is a wrong
implementation even if it matches.

%=====================================================================
\newpage
\problemheader{Problem A3: Circular Convolution and Wraparound}{A1, A2}{a3\_circular\_wraparound.py}

A length-$N$ transform never computes linear convolution. It computes
\emph{circular} convolution of period $N$: anything that would land past slot
$N-1$ is added back onto slot $m \bmod N$. The padding rule exists only to make
the two coincide. In this problem you compute the wrong thing on purpose and
prove exactly how it is wrong.

You are given two reproducible operands in \file{a3\_circular\_wraparound.py}.
Complete only the three blocks marked ``TODO''. Do not modify any offline file.

\textbf{Mathematical model.} With linear convolution $\ell = a * b$ of length
$n + q - 1$ and any $N \ge \max(n, q)$,
\[
  (a \circledast_N b)[m] \;=\; \sum_{k \,\equiv\, m \ (\mathrm{mod}\ N)} \ell[k],
  \qquad m = 0, \dots, N-1 .
\]
The left side is what a length-$N$ transform returns; the right side is the
linear result \emph{folded} onto $N$ slots. They must agree exactly, integer
for integer. The offline's \code{multiply\_transform} supplies $\ell$.

\textbf{Tasks}
\begin{enumerate}[nosep]
  \item Complete \code{short\_length}. Return $\max(n, q)$ rounded to what the
        engine supports --- enough room for either input, deliberately too
        little for the product.
        \hfill [2 Marks]
  \item Complete \code{circular\_convolve}. Pad both arrays to exactly $N$,
        transform, multiply, inverse-transform, keep the real part, round.
        Return all $N$ slots.
        \hfill [4 Marks]
  \item Complete \code{fold}. Wrap a linear result onto $N$ slots:
        \code{out[m \% N] += linear[m]}.
        \hfill [4 Marks]
\end{enumerate}

\textbf{Running and verification.}
\begin{lstlisting}
python a3_circular_wraparound.py --digits 120 --seed 5 --engine fft
\end{lstlisting}
The runner checks $\max_m |\text{circular}[m] - \text{folded}[m]| = 0$ and that
the padded route still equals \code{int(a) * int(b)}. The report also prints
the plausible-looking wrong number the short transform produces, next to the
correct one. It must show \code{verification: MATCH}.

%=====================================================================
\newpage
\problemheader{Problem A4: Shift Detection by Cross-Correlation}{A1, A2}{a4\_shift\_correlation.py}

Correlation slides one sequence past another and measures how well they line
up at each offset. It is convolution with a time-reversed sequence, and
reversing a real sequence conjugates its spectrum --- so correlation costs the
same three transforms as multiplication, with one conjugate.

You are given a number $a$ and $b = a \cdot B^{s}$ (the same limbs moved up by
$s$ positions) in \file{a4\_shift\_correlation.py}. Complete only the three
blocks marked ``TODO''. Do not modify any offline file.

\textbf{Mathematical model.} For limb arrays $a$ (length $n$) and $b$ (length
$q$) the cross-correlation at lag $k$ is
\[
  c[k] = \sum_i a[i]\, b[i+k], \qquad
  \mathrm{DFT}\{c\} = \overline{A} \cdot B, \qquad N \ge n + q - 1,
\]
with lag $k$ stored at index $k \bmod N$ (negative lags sit at the end of the
array). Because $b$ is $a$ shifted by $s$, the largest correlation occurs at
$k = s$ and equals $\sum_i a[i]^2$, the energy of $a$.

\textbf{Tasks}
\begin{enumerate}[nosep]
  \item Complete \code{choose\_transform\_length}. Every lag needs its own
        slot: the linear-convolution rule applies unchanged.
        \hfill [2 Marks]
  \item Complete \code{cross\_correlate}. Pad, transform both, multiply
        $\overline{A}$ by $B$, perform one inverse transform, keep the real
        part, round to \code{int64}. Return the full length-$N$ array and $N$.
        \hfill [5 Marks]
  \item Complete \code{find\_shift}. Correlate the limb magnitudes, locate the
        maximum, unwrap an index above $N/2$ to a negative lag, return the lag,
        the peak value and $N$.
        \hfill [3 Marks]
\end{enumerate}

\textbf{Running and verification.}
\begin{lstlisting}
python a4_shift_correlation.py --digits 400 --shift 7 --engine fft
\end{lstlisting}
The runner recomputes every lag with the defining double loop, checks the
detected shift and checks that the peak equals the energy of $a$. The report
must show \code{verification: MATCH}.

%=====================================================================
\newpage
\problemheader{Problem A5: Fourth Power by Repeated Squaring}{A1, A2}{a5\_fourth\_power.py}

Squaring is multiplication with one operand, so it needs only \emph{one}
forward transform: $A \cdot A$. A fourth power is a square of a square. The
catch is between the two squarings: the coefficients that come back from the
first are not digits yet, and squaring un-carried coefficients breaks the
mantissa bound.

You are given the first operand of an offline input file in
\file{a5\_fourth\_power.py}. Complete only the three blocks marked ``TODO''. Do
not modify any offline file.

\textbf{Mathematical model.} For an $n$-limb array $a$ with spectrum $A$,
\[
  a^2 = \mathrm{IDFT}\{A \cdot A\} \quad (2n - 1 \text{ coefficients}),
  \qquad a^4 = \big(\,\mathrm{carry}(a^2)\,\big)^2 .
\]
The un-carried coefficients of $a^2$ reach $n(B-1)^2 \approx 10^{11}$; squaring
them directly would produce sums near $10^{25}$, far past $2^{53}$. Carrying
first (\code{from\_limbs}, then \code{to\_limbs}) restores limbs below $B$, and
the second squaring is safe. An even power is never negative.

\textbf{Tasks}
\begin{enumerate}[nosep]
  \item Complete \code{square\_length}. A square has $2n - 1$ coefficients;
        round to what the engine supports.
        \hfill [2 Marks]
  \item Complete \code{square\_transform}. One forward transform, square the
        spectrum, one inverse, real part, crop, round. Two forward transforms
        is wrong.
        \hfill [4 Marks]
  \item Complete \code{fourth\_power}. Square, carry into a proper number,
        re-split into limbs, square again, carry, return the string together
        with both transform lengths and both limb counts.
        \hfill [4 Marks]
\end{enumerate}

\textbf{Running and verification.}
\begin{lstlisting}
python a5_fourth_power.py --input inputs/3.txt --engine fft
\end{lstlisting}
The runner compares against \code{int(a) ** 4}. The report must show
\code{verification: MATCH}.

%=====================================================================
\newpage
\problemheader{Problem B1: Shifting an Image With a Phase Ramp}{B1, B2}{b1\_spectral\_shift.py}

Delaying a signal does not change how much of each frequency it contains, only
where each oscillation sits --- its phase. The shift theorem makes that exact:
a delay of $s$ samples multiplies the spectrum by a linear phase ramp. In two
dimensions the ramps for rows and columns simply add in the exponent.

You are given an image and a shift in \file{b1\_spectral\_shift.py}. Complete
only the three blocks marked ``TODO''. Do not modify \file{transforms.py},
\file{image\_conv.py}, or any other offline file.

\textbf{Mathematical model.} For an $H \times W$ plane with spectrum $F[u,v]$
and integer shifts $(d_r, d_c)$,
\[
  F_{\text{shifted}}[u,v] \;=\; F[u,v]\; e^{-2\pi j \left( \frac{u\, d_r}{H} + \frac{v\, d_c}{W} \right)},
  \qquad u = 0..H-1,\ v = 0..W-1 .
\]
No padding is used, so the shift wraps around the edges: the result is a
\emph{circular} shift, which is exactly what \code{np.roll} computes. Negative
shifts are allowed.

\textbf{Tasks}
\begin{enumerate}[nosep]
  \item Complete \code{phase\_ramp}. Build the $H \times W$ complex multiplier
        above from the two index vectors.
        \hfill [3 Marks]
  \item Complete \code{shift\_plane}. Transform the plane, multiply by the
        ramp, perform exactly one inverse 2D transform, keep the real part.
        \hfill [4 Marks]
  \item Complete \code{shift\_image}. One plane for grayscale, three planes
        independently for RGB, original shape preserved.
        \hfill [3 Marks]
\end{enumerate}

\textbf{Running and verification.}
\begin{lstlisting}
python b1_spectral_shift.py --image images/skyline256.png --rows 37 --cols -21 --engine fft
python b1_spectral_shift.py --color
\end{lstlisting}
The runner compares against \code{np.roll(image, (rows, cols), axis=(0, 1))}
and writes \file{shifted.png}, \file{comparison.png} and \file{report.txt}. The
report must show \code{verification: MATCH}. Both grayscale and RGB inputs must
work.

%=====================================================================
\newpage
\problemheader{Problem B2: Unsharp Masking in the Frequency Domain}{B1, B2}{b2\_unsharp\_mask.py}

A blur removes detail; subtracting the blur from the original isolates that
detail; adding some of it back sharpens the picture. That is unsharp masking,
and every step of it is a convolution or a sum, so the whole thing is one
expression in the frequency domain.

You are given an image and a Gaussian kernel in \file{b2\_unsharp\_mask.py}.
Complete only the three blocks marked ``TODO''. Do not modify any offline file.

\textbf{Mathematical model.} With the padded spectra $P$ (image) and $G$
(kernel), the centred impulse spectrum $\Delta$ (provided helper) and a
sharpening amount $\lambda$,
\[
  y = x + \lambda\,(x - x * g)
  \quad\Longleftrightarrow\quad
  Y = P\,\big(\Delta + \lambda\,(\Delta - G)\big) = P\,\big((1+\lambda)\Delta - \lambda G\big).
\]
Use one inverse transform and the usual linear-convolution crop starting at
$(\lfloor K_r/2 \rfloor, \lfloor K_c/2 \rfloor)$. An $R \times C$ image and a
$K_r \times K_c$ kernel need at least $(R + K_r - 1) \times (C + K_c - 1)$
samples; the radix-2 FFT needs each dimension raised to the next power of two.

\textbf{Tasks}
\begin{enumerate}[nosep]
  \item Complete \code{choose\_transform\_shape}. Minimum linear-convolution
        dimensions, then the engine's restriction.
        \hfill [2 Marks]
  \item Complete \code{sharpen\_plane}. Pad, transform, form $Y$, perform
        exactly one inverse 2D transform, keep the real part, crop.
        \hfill [5 Marks]
  \item Complete \code{sharpen\_image}. Grayscale or RGB dispatch, shape
        preserved.
        \hfill [3 Marks]
\end{enumerate}

\textbf{Running and verification.}
\begin{lstlisting}
python b2_unsharp_mask.py --image images/skyline256.png --kernel-size 21 --amount 1.5 --engine fft
python b2_unsharp_mask.py --color
\end{lstlisting}
The runner computes \code{image + amount * (image - convolve\_image(image, kernel))}
as an independent oracle. The report must show \code{verification: MATCH}.

%=====================================================================
\newpage
\problemheader{Problem B3: Deblurring by Inverse Filtering}{B1, B2}{b3\_deblur.py}

If a blur is a pointwise multiplication in the frequency domain, undoing it is
a pointwise division --- as long as nothing is divided by zero. A circular blur
by a small Gaussian can be undone exactly; the same idea applied to a large
kernel amplifies rounding noise into garbage, and you will see why.

You are given a circularly blurred image and the kernel that blurred it in
\file{b3\_deblur.py}. Complete only the three blocks marked ``TODO''. Do not
modify any offline file.

\textbf{Mathematical model.} The offline's circular path computes
$Y = X \cdot G_w$, where $G_w$ is the spectrum of the kernel wrapped around the
origin (centre tap at index $(0,0)$). Then
\[
  X[u,v] =
  \begin{cases}
    Y[u,v] / G_w[u,v], & |G_w[u,v]| > \varepsilon, \\
    0, & \text{otherwise},
  \end{cases}
  \qquad x = \mathrm{Re}\,\mathrm{IDFT}\{X\}.
\]
Rounding noise of size $10^{-16}$ in $Y$ becomes $10^{-16} / |G_w|$ in $X$: a
$5 \times 5$ Gaussian keeps $|G_w| \ge 4 \times 10^{-3}$ and restores the image
to $10^{-14}$; a $31 \times 31$ Gaussian has bins near $10^{-50}$.

\textbf{Tasks}
\begin{enumerate}[nosep]
  \item Complete \code{wrapped\_kernel\_spectrum}. Embed the kernel at the
        origin of a zero array of the plane's shape, roll by
        $(-\lfloor K_r/2 \rfloor, -\lfloor K_c/2 \rfloor)$, transform.
        \hfill [3 Marks]
  \item Complete \code{deblur\_plane}. Divide only where $|G_w| > \varepsilon$,
        zero elsewhere, one inverse transform, real part. Never divide by a
        vanishing bin.
        \hfill [4 Marks]
  \item Complete \code{deblur\_image}. Grayscale or RGB dispatch, shape
        preserved.
        \hfill [3 Marks]
\end{enumerate}

\textbf{Running and verification.}
\begin{lstlisting}
python b3_deblur.py --image images/skyline256.png --kernel-size 5 --engine fft
python b3_deblur.py --color
\end{lstlisting}
The runner blurs with \code{convolve\_image(..., circular=True)}, calls your
deblur and compares with the original image. The report must show
\code{verification: MATCH}. Afterwards try \code{--kernel-size 31} and read the
``smallest $|G|$ bin'' line of the report.

%=====================================================================
\newpage
\problemheader{Problem B4: Phase Correlation}{B1, B2}{b4\_phase\_correlation.py}

Online B showed that phase carries the structure of an image. Phase
correlation uses that: if two images differ only by a shift, their spectra
differ only by a phase ramp, and dividing that ramp out of the cross-power
spectrum leaves a pure impulse whose position \emph{is} the shift.

You are given an image and its shifted copy in
\file{b4\_phase\_correlation.py}. Complete only the three blocks marked
``TODO''. Do not modify any offline file.

\textbf{Mathematical model.} For planes $a$ and $b = \mathrm{roll}(a, (d_r, d_c))$
with spectra $A$ and $B$,
\[
  R[u,v] = \frac{\overline{A[u,v]}\, B[u,v]}{\big|\overline{A[u,v]}\, B[u,v]\big|}
  \;=\; e^{-2\pi j \left( \frac{u\, d_r}{H} + \frac{v\, d_c}{W} \right)},
  \qquad r = \mathrm{Re}\,\mathrm{IDFT}\{R\} = \delta[m - d_r,\; n - d_c].
\]
Bins with $|\overline{A} B| \le \varepsilon$ must be set to $0$, never divided.
The peak of $r$ at $(m, n)$ gives the shift; an index above $H/2$ (or $W/2$)
represents a negative shift $m - H$ (or $n - W$). For RGB, sum the three
channel surfaces before locating the peak.

\textbf{Tasks}
\begin{enumerate}[nosep]
  \item Complete \code{unit\_cross\_power}. Form $\overline{A} B$, normalise
        reliable bins to unit magnitude, zero the rest, no division by zero.
        \hfill [3 Marks]
  \item Complete \code{correlation\_surface}. Transform both planes, take the
        unit cross-power spectrum, one inverse transform, real part.
        \hfill [4 Marks]
  \item Complete \code{find\_shift}. Grayscale or RGB dispatch, locate the
        peak, unwrap it, return \code{(rows, cols, surface)}.
        \hfill [3 Marks]
\end{enumerate}

\textbf{Running and verification.}
\begin{lstlisting}
python b4_phase_correlation.py --image images/skyline256.png --rows 37 --cols -21 --engine fft
python b4_phase_correlation.py --color
\end{lstlisting}
The runner shifts with \code{np.roll}, checks the detected shift and checks
that the surface is a unit impulse to $10^{-9}$. The report must show
\code{verification: MATCH}.

%=====================================================================
\newpage
\problemheader{Problem B5: Ideal Low-Pass / High-Pass Split}{B1, B2}{b5\_ideal\_filter.py}

Low frequencies describe broad, slow variation; high frequencies describe
edges and texture. An ideal filter keeps every bin inside a radius of zero
frequency and discards the rest --- but in a DFT array zero frequency is at
index $0$ and the negative frequencies live at the \emph{top} of the array, so
the distance must be measured with wrapped indices.

You are given an image and a radius in \file{b5\_ideal\_filter.py}. Complete
only the three blocks marked ``TODO''. Do not modify any offline file.

\textbf{Mathematical model.} Bin $u$ of an $H$-point DFT has frequency
$f_u = \min(u, H - u)$, likewise $f_v = \min(v, W - v)$. With
\[
  M[u,v] = \begin{cases} 1, & \sqrt{f_u^2 + f_v^2} \le r \\ 0, & \text{otherwise} \end{cases},
  \qquad
  x_{\text{low}} = \mathrm{Re}\,\mathrm{IDFT}\{F M\}, \quad
  x_{\text{high}} = \mathrm{Re}\,\mathrm{IDFT}\{F (1 - M)\},
\]
linearity gives $x_{\text{low}} + x_{\text{high}} = x$ exactly, and Parseval's
theorem says the fraction of energy kept is the same whether measured on pixels
or on spectrum bins.

\textbf{Tasks}
\begin{enumerate}[nosep]
  \item Complete \code{radial\_mask}. Wrapped frequency distance on both axes,
        Euclidean radius test, return \code{float64}.
        \hfill [3 Marks]
  \item Complete \code{split\_plane}. One forward transform, two masked
        inverse transforms, real parts, return \code{(low, high)}.
        \hfill [4 Marks]
  \item Complete \code{split\_image}. Grayscale returns the pair directly; RGB
        splits each channel, stacks the three lows and the three highs
        separately, returns both.
        \hfill [3 Marks]
\end{enumerate}

\textbf{Running and verification.}
\begin{lstlisting}
python b5_ideal_filter.py --image images/skyline256.png --radius 20 --engine fft
python b5_ideal_filter.py --color
\end{lstlisting}
The runner checks $\max|x_{\text{low}} + x_{\text{high}} - x|$, compares your
low-pass with one built from an independently constructed (centred, then
rolled) mask, and checks Parseval. The report must show
\code{verification: MATCH}.

%=====================================================================
\newpage
\section*{Solutions in plain language}

Every problem here is the same five-step pipeline you already wrote in the
offline: \emph{pad} $\to$ \emph{transform} $\to$ \emph{do something simple to
the spectrum} $\to$ \emph{inverse transform once} $\to$ \emph{real part, crop,
round}. What changes is only the ``something simple''. Full code is in
\file{practice/solutions/}; below is the idea behind each, written for a
first-time reader.

\subsection*{The three facts everything rests on}

\begin{enumerate}[nosep]
  \item \textbf{Convolution becomes multiplication.} The DFT rewrites a signal
        as a sum of complex exponentials. Convolution scales each exponential
        and never mixes it with its neighbours, so after the transform the
        whole tangled sum turns into $N$ independent multiplications. Hence
        $a * b \leftrightarrow A \cdot B$, and with more factors
        $a * b * c \leftrightarrow A \cdot B \cdot C$.
  \item \textbf{The DFT is linear.} $\mathrm{DFT}\{x + y\} = X + Y$ and
        $\mathrm{DFT}\{\lambda x\} = \lambda X$. So any expression built from
        sums and convolutions can be assembled entirely in the frequency domain
        and brought back with a single inverse transform.
  \item \textbf{A shift is a phase ramp.} Delaying $x[n]$ by $s$ samples turns
        $X[u]$ into $X[u]\,e^{-2\pi j u s / N}$: same magnitudes, rotated
        phases. Reversing a real sequence in time conjugates its spectrum.
\end{enumerate}
And one rule of engineering: a length-$N$ transform computes \emph{circular}
convolution, so pad to at least the linear length or the tail wraps onto the
head.

\subsection*{A1 --- product of three integers}

\textbf{Idea.} Fact~1 with three factors. Transform $a$, $b$, $c$ once each,
multiply the three spectra bin by bin, invert once.

\textbf{Why the base drops to $10^2$.} Every coefficient of the result comes
back as a floating-point number that you round. Rounding only recovers the
correct integer if the integer was representable: doubles hold integers
exactly up to $2^{53} \approx 9 \times 10^{15}$. A coefficient of $a b c$ adds
up about $n^2/2$ terms, each a product of three limbs, so it can reach
$(n^2/2)(B-1)^3$. With $B = 10^4$ and a thousand limbs that is $10^{18}$ ---
too big. With $B = 10^2$ it is $10^{12}$ --- fine. This is the same argument
the offline made for choosing $10^4$ over $10^9$, one factor further.

\textbf{Steps.} TODO~1: \code{need = n + q + r - 2}, power of two unless the
engine is \code{arbitrary}. TODO~2: three zero-padded arrays, three
\code{engine.transform}, one product, one \code{engine.inverse}, \code{.real},
crop to \code{need}, \code{np.rint}. TODO~3: \code{to\_limbs(text, 2)} three
times, product of the three signs, \code{from\_limbs(sign, coeffs, 2)}.

\textbf{Pitfall.} Forgetting to pass the base to \code{from\_limbs}: the
carry sweep then uses $10^4$ on base-$10^2$ limbs and prints a wrong number
with the right length.

\subsection*{A2 --- $ab + cd$ with one inverse}

\textbf{Idea.} Fact~2. $\mathrm{DFT}\{a*b + c*d\} = AB + CD$. Four forward
transforms, one addition, one inverse. Doing two inverses and adding the
results also gives the right number, but it misses the point of the problem
and is marked as such.

\textbf{Why one $N$ is enough.} Both products must fit without wrapping, so
$N$ must be at least the longer of the two linear lengths; padding the shorter
product with extra zeros costs nothing.

\textbf{Why non-negative only.} $ab$ and $cd$ are both non-negative, so their
sum's coefficients are non-negative and the carry sweep works. If one product
were negative, some coefficients would be negative and \code{from\_limbs}
would loop forever on a negative carry --- hence the \code{ValueError}.

\textbf{Steps.} TODO~1: \code{max(la+lb-1, lc+ld-1)}, engine rounding.
TODO~2: pad four, transform four, \code{A*B + C*D}, one inverse.
TODO~3: convert, check signs, convolve, \code{from\_limbs(1, coeffs)}.

\subsection*{A3 --- circular convolution and wraparound}

\textbf{Idea.} The transform of length $N$ ``thinks'' the world is a circle of
$N$ slots. A coefficient that should land in slot $m \ge N$ lands in slot
$m \bmod N$ instead, on top of whatever was already there. That is not an
approximation --- it is exact arithmetic, just of the wrong quantity. The
problem makes you compute both sides of the identity and show they are equal.

\textbf{Steps.} TODO~1: \code{max(la, lb)}, engine rounding --- a length that
is big enough for the inputs but not for the product. TODO~2: identical to
\code{multiply\_transform} except that $N$ is given and nothing is cropped.
TODO~3: a loop adding \code{linear[m]} into \code{out[m \% N]}.

\textbf{What to notice in the report.} The wrapped ``product'' has fewer digits
and a completely different start; nothing crashed and no warning appeared. This
is the silent failure the padding rule prevents.

\subsection*{A4 --- shift detection by cross-correlation}

\textbf{Idea.} Correlation $c[k] = \sum_i a[i]\,b[i+k]$ asks ``how well does
$b$ match $a$ when slid by $k$?'' It is convolution with $a$ reversed. Fact~3:
reversing a real sequence conjugates its spectrum, so
$\mathrm{DFT}\{c\} = \overline{A}\,B$. One \code{np.conj} is the entire
difference from multiplication.

\textbf{Why the peak is at $s$.} Appending $4s$ decimal zeros multiplies by
$10^{4s} = B^{s}$, which moves every limb up $s$ places. Sliding by exactly $s$
lines each limb up with itself, and the sum of squares is as large as it can
be. Every other lag mixes different limbs and comes out smaller.

\textbf{Where negative lags live.} There are $n + q - 1$ possible lags. Lag
$k \ge 0$ sits at index $k$; lag $k < 0$ sits at index $N + k$. If the peak
index is above $N/2$ it is really a negative lag, so subtract $N$.

\textbf{Steps.} TODO~1: \code{la + lb - 1}, engine rounding. TODO~2: pad,
\code{np.conj(engine.transform(fa)) * engine.transform(fb)}, inverse,
\code{.real}, \code{np.rint}. TODO~3: \code{np.argmax}, unwrap, return.

\subsection*{A5 --- fourth power by repeated squaring}

\textbf{Idea.} $a^2 = \mathrm{IDFT}\{A \cdot A\}$ needs one forward transform,
not two. $a^4 = (a^2)^2$.

\textbf{Why carry in between.} After the first squaring, the coefficients are
sums like $\sum a[i] a[m-i]$, up to $n (B-1)^2 \approx 10^{11}$ each --- far
larger than a limb. They are a correct answer, just not written in base $B$.
Squaring them again would produce coefficients around $10^{25}$, past the
mantissa, and rounding would return wrong integers. Running the carry sweep
(\code{from\_limbs}) and re-splitting (\code{to\_limbs}) rewrites $a^2$ as
proper limbs below $B$, and then the second squaring is exactly as safe as the
first.

\textbf{Steps.} TODO~1: \code{2n - 1}, engine rounding. TODO~2: one
\code{engine.transform}, \code{A*A}, inverse, crop, round. TODO~3:
\code{square}, \code{to\_limbs(from\_limbs(1, square))}, \code{square} again,
\code{from\_limbs}.

\subsection*{B1 --- shifting an image with a phase ramp}

\textbf{Idea.} Fact~3 in two dimensions. Moving the picture down by $d_r$ rows
and right by $d_c$ columns multiplies bin $(u, v)$ by
$e^{-2\pi j (u d_r / H + v d_c / W)}$. Magnitudes are untouched --- the picture
contains the same frequencies --- only phases turn.

\textbf{Why it wraps.} The DFT treats the image as periodic, so content pushed
off the bottom edge re-enters at the top. That is exactly \code{np.roll}, which
is why it is the oracle. No padding, no crop.

\textbf{Steps.} TODO~1: two index vectors \code{u[:, None]} and
\code{v[None, :]}, one \code{np.exp}. TODO~2: \code{transform\_2d}, multiply,
\code{inverse\_2d}, \code{.real}. TODO~3: the standard dispatch.

\textbf{Pitfall.} A plus sign in the exponent shifts the other way; the
oracle catches it immediately.

\subsection*{B2 --- unsharp masking}

\textbf{Idea.} Detail $= x - x * g$ (original minus blur). Sharpened
$= x + \lambda \cdot \text{detail}$. In the frequency domain, $x * g$ is $P G$
and the unblurred $x$ is $P \Delta$ --- $\Delta$ is the spectrum of an impulse
sitting at the kernel's centre, so that after the same crop it lands exactly on
the original pixels (that is the whole reason Online A provided it). Then
$Y = P\Delta + \lambda(P\Delta - PG) = P(\Delta + \lambda(\Delta - G))$, one
inverse transform, one crop.

\textbf{Steps.} TODO~1: exactly Online A's \code{choose\_transform\_shape}.
TODO~2: pad plane and kernel, two \code{transform\_2d}, the formula, one
\code{inverse\_2d}, \code{.real}, crop at \code{(kh//2, kw//2)}. TODO~3: the
standard dispatch.

\textbf{Pitfall.} Using $1$ instead of $\Delta$ for ``the original'': a plain
$1$ is an impulse at the corner, not at the kernel centre, so the original
comes out displaced by half a kernel relative to the blur.

\subsection*{B3 --- deblurring by inverse filtering}

\textbf{Idea.} A circular blur is $Y = X G_w$ bin by bin, so $X = Y / G_w$ bin
by bin. Division is only allowed where $G_w$ is not (numerically) zero.

\textbf{Why the kernel must be wrapped.} The offline's circular path puts the
kernel's centre tap at index $(0,0)$ with \code{np.roll}. The blur you are
undoing was made that way, so the spectrum you divide by must be of the
identically placed kernel, or every pixel comes back shifted by half a kernel.

\textbf{Why big kernels cannot be inverted.} A blur kills high frequencies:
$|G_w|$ is tiny there. Dividing by a tiny number multiplies whatever is in $Y$
at that bin --- including $10^{-16}$ rounding noise --- by a huge factor. For a
$5 \times 5$ Gaussian the smallest bin is about $4 \times 10^{-3}$, so noise
grows $250\times$ and stays invisible. For $31 \times 31$ the smallest bins are
around $10^{-50}$; the $\varepsilon$ mask zeros those, but bins just above
$\varepsilon$ still amplify by $10^{12}$. This is the mantissa argument again,
from the other side.

\textbf{Steps.} TODO~1: zeros of the plane's shape, kernel at top-left,
\code{np.roll} by \code{(-(kh//2), -(kw//2))}, \code{transform\_2d}. TODO~2:
mask $|G_w| > \varepsilon$, divide at the mask only, zeros elsewhere, one
inverse, \code{.real}. TODO~3: the standard dispatch.

\subsection*{B4 --- phase correlation}

\textbf{Idea.} If $b$ is $a$ shifted, then by Fact~3 $B = A \cdot \text{ramp}$.
Multiply by $\overline{A}$: $\overline{A} B = |A|^2 \cdot \text{ramp}$. Divide
by the magnitude: only the ramp is left. The inverse transform of a pure ramp
is a single impulse of height $1$ at the shift --- an entire image collapses to
one bright pixel whose position is the answer.

\textbf{Why normalise.} Without the division the surface is
$\mathrm{IDFT}\{|A|^2 \text{ramp}\}$ = the autocorrelation of $a$ shifted ---
a broad blob, still peaked at the right place but far less sharp. Keeping only
phase is what makes the peak an impulse; that is the same lesson as Online B.

\textbf{Steps.} TODO~1: the \code{unit\_phase} pattern from Online B applied to
\code{np.conj(A) * B}, with zeros (not ones) for unreliable bins. TODO~2: two
\code{transform\_2d}, TODO~1, one \code{inverse\_2d}, \code{.real}. TODO~3:
sum the channel surfaces for RGB, \code{np.argmax} + \code{np.unravel\_index},
subtract $H$ or $W$ from an index above the half.

\textbf{Pitfall.} $A\,\overline{B}$ instead of $\overline{A}\,B$ puts the peak
at $(-d_r, -d_c)$.

\subsection*{B5 --- ideal low-pass / high-pass split}

\textbf{Idea.} Multiply the spectrum by a mask that is $1$ near zero frequency
and $0$ far from it: low-pass. The complementary mask $1 - M$: high-pass. By
Fact~2, $F M + F (1 - M) = F$, so the two parts add back to the original
exactly.

\textbf{Why wrapped distances.} In a DFT array bin $0$ is zero frequency, bin
$1$ is the slowest oscillation, and bin $H-1$ is \emph{also} the slowest
oscillation (going the other way). The frequency of bin $u$ is therefore
$\min(u, H - u)$, not $u$. Measuring plain distance from index $(0,0)$ would
treat the top-right corner as high frequency when it is low.

\textbf{Why Parseval holds.} The DFT is an orthogonal change of basis (up to
the factor $N$), so it preserves lengths: $\sum |x|^2 = \frac{1}{N} \sum |X|^2$.
The fraction of energy inside the mask is the same whether you sum pixels of
$x_{\text{low}}$ or bins of $F M$; the runner checks that to $10^{-9}$.

\textbf{Steps.} TODO~1: \code{np.minimum(u, H - u)}, same for columns, outer
Euclidean distance, \code{<= radius}, \code{.astype(float64)}. TODO~2: one
\code{transform\_2d}, two \code{inverse\_2d}, \code{.real} on both. TODO~3:
the pair-returning dispatch --- collect \code{(low, high)} per channel, then
\code{np.stack} the lows and the highs separately.

\textbf{Pitfall.} Returning a bool mask: \code{spectrum * mask} with a bool
array works, but \code{1 - mask} on a bool array does not mean what you think.
Cast to \code{float64}.

\subsection*{The dispatch you will write every single time}

\begin{lstlisting}
if image.ndim == 2:
    return f(image, ...)
if image.ndim == 3 and image.shape[2] == 3:
    planes = [f(image[:, :, c], ...) for c in range(image.shape[2])]
    return np.stack(planes, axis=-1)
raise ValueError("images must be grayscale or RGB")
\end{lstlisting}
When \code{f} returns a pair, replace the last two lines of the RGB branch by
\code{firsts, seconds = zip(*pairs)} and stack each with \code{axis=-1}. Type it
from memory before the online starts.

\end{document}
